\chapter{pg notes}

\begin{enumerate}
        \item Int\'egration d'outils
        \begin{itemize}
                \item interface humaine (X) param\'etreurs et d\'eclencheurs
                \item outils de visualisation
                \item outils d'analyse: transformatons d'espace
                \item outils de???
        \end{itemize}
        \item anneal
        \begin{itemize}
                \item temps
                \item densit\'ee
                \item volume
        \end{itemize}
        \item m\'ethodes num\'eriques\\
                suggestion PG: voir {em transform\'ees}
                Fourrier, Walsh, Wavelet, Hartley.
                Automates de Fourrier.\\
                suggestion PG:
                \begin{enumerate}
                        \item faire de la recherche assist\'e par AC.
                        \item id\'ee d'``automates de phas\'e' a pixels de phases.
                \end{enumerate}
        \item misc
        \begin{itemize}
                \item collecte de voisinnage.
                \item combinatoire de r\`egles.
                \item compilateur de voisinnage et de table (LUT)
        \end{itemize}
        \item Probl\`eme de la relation entre la r\`egle et son bruitage
                et l'espace initial.\\
                Et critique des notions de fronti\`ere du chaos.\\
                PG: fronti\`ere d'universalit\'e.\\
                PG: universalit\'e: ``gestion de table glorifi\'e''
        \item suggestion PG:
        \begin{enumerate}
                \item plac\'e la dedans des crabes.
                \item les crabes sont de automates avec des voisinnage
                        de tables et des macro-m\'eta r\`egles
                        (\'equivalent \`a des syst\`emes experts dot\'es
                        de connaissances sur les cas).
                \item diagonaliser lees crabes.
        \end{enumerate}
\end{enumerate}

%%
%% Local Variables:
%% mode: auto-fill
%% iso-accents-minor-mode: nil
%% TeX-master: "top"
%% TeX-command-default: "mytex"
%% Local IspellParsing: "latex-mode ~latin1"
%% Local IspellPersDict: "~/.ispell_lisp"
%% End:
